\documentclass[sigconf]{acmart}

\usepackage{polyglossia}
\setmainlanguage{vietnamese}
\setotherlanguage{english}

\usepackage{pifont}
\usepackage[capitalise,english]{cleveref}

% Placeholder cho số liệu chưa có: các thí nghiệm trên dải Qwen đang chạy;
% số cuối cùng được điền qua generated/*.tex do pipeline thống kê sinh tự động.
\newcommand{\TBD}[1]{\textcolor{red}{[TBD: #1]}}

% Macro thống kê sinh tự động từ artifact (demo placeholder khi chưa có số).
% AUTO-GENERATED by isd-agent-benchmark/scripts/8_gen_paper_tables.py — DO NOT EDIT BY HAND
% stats_macros.tex — headline numbers for prose
% generated_at: 2026-07-21T02:19:17
% !!!!!!!!!!!!!!!!!!!!!!!!!!!!!!!!!!!!!!!!!!!!!!!!!!!!!!!!!!!!!!!!!!!!!!
% !! TBD-DEMO: ALL VALUES BELOW ARE FAKE PLACEHOLDERS (--demo mode). !!
% !! Regenerate from the real pooled_ladder.json before submission.  !!
% !!!!!!!!!!!!!!!!!!!!!!!!!!!!!!!!!!!!!!!!!!!!!!!!!!!!!!!!!!!!!!!!!!!!!!

% number of model sizes in the ladder
\newcommand{\ladderNumModels}{4}
% independent runs per model size
\newcommand{\ladderNumRunsPerModel}{3}
% scenarios per run (test_90)
\newcommand{\ladderNumScenarios}{90}
% agents compared (6 baselines + proposed)
\newcommand{\ladderNumAgents}{7}
% alignmentgraph-isd pooled mean Total at qwen3.5-0.8b
\newcommand{\rqOneHarnessTotalZeroEightB}{46{,}00}
% alignmentgraph-isd run-to-run SD (SD of the run-level means) at qwen3.5-0.8b
\newcommand{\rqOneHarnessRunSDZeroEightB}{0{,}40}
% best baseline (dick-carey-agent) pooled mean Total at qwen3.5-0.8b
\newcommand{\rqOneBestBaselineTotalZeroEightB}{42{,}00}
% name of the best baseline at qwen3.5-0.8b
\newcommand{\rqOneBestBaselineNameZeroEightB}{Dick-Carey-Agent}
% mean paired delta alignmentgraph-isd - dick-carey-agent at qwen3.5-0.8b
\newcommand{\rqOneDeltaBestZeroEightB}{4{,}00}
% bootstrap 95% CI lower bound of that delta
\newcommand{\rqOneDeltaBestCILoZeroEightB}{3{,}00}
% bootstrap 95% CI upper bound of that delta
\newcommand{\rqOneDeltaBestCIHiZeroEightB}{5{,}00}
% max Holm-adjusted p over the 6 comparisons at qwen3.5-0.8b
\newcommand{\rqOnePHolmMaxZeroEightB}{6{,}0\times 10^{-9}}
% alignmentgraph-isd pooled mean Total at qwen3.5-2b
\newcommand{\rqOneHarnessTotalTwoB}{61{,}00}
% alignmentgraph-isd run-to-run SD (SD of the run-level means) at qwen3.5-2b
\newcommand{\rqOneHarnessRunSDTwoB}{0{,}40}
% best baseline (dick-carey-agent) pooled mean Total at qwen3.5-2b
\newcommand{\rqOneBestBaselineTotalTwoB}{57{,}00}
% name of the best baseline at qwen3.5-2b
\newcommand{\rqOneBestBaselineNameTwoB}{Dick-Carey-Agent}
% mean paired delta alignmentgraph-isd - dick-carey-agent at qwen3.5-2b
\newcommand{\rqOneDeltaBestTwoB}{4{,}00}
% bootstrap 95% CI lower bound of that delta
\newcommand{\rqOneDeltaBestCILoTwoB}{3{,}00}
% bootstrap 95% CI upper bound of that delta
\newcommand{\rqOneDeltaBestCIHiTwoB}{5{,}00}
% max Holm-adjusted p over the 6 comparisons at qwen3.5-2b
\newcommand{\rqOnePHolmMaxTwoB}{6{,}0\times 10^{-9}}
% alignmentgraph-isd pooled mean Total at qwen3.5-4b
\newcommand{\rqOneHarnessTotalFourB}{71{,}00}
% alignmentgraph-isd run-to-run SD (SD of the run-level means) at qwen3.5-4b
\newcommand{\rqOneHarnessRunSDFourB}{0{,}40}
% best baseline (dick-carey-agent) pooled mean Total at qwen3.5-4b
\newcommand{\rqOneBestBaselineTotalFourB}{67{,}00}
% name of the best baseline at qwen3.5-4b
\newcommand{\rqOneBestBaselineNameFourB}{Dick-Carey-Agent}
% mean paired delta alignmentgraph-isd - dick-carey-agent at qwen3.5-4b
\newcommand{\rqOneDeltaBestFourB}{4{,}00}
% bootstrap 95% CI lower bound of that delta
\newcommand{\rqOneDeltaBestCILoFourB}{3{,}00}
% bootstrap 95% CI upper bound of that delta
\newcommand{\rqOneDeltaBestCIHiFourB}{5{,}00}
% max Holm-adjusted p over the 6 comparisons at qwen3.5-4b
\newcommand{\rqOnePHolmMaxFourB}{6{,}0\times 10^{-9}}
% alignmentgraph-isd pooled mean Total at qwen3.5-9b
\newcommand{\rqOneHarnessTotalNineB}{81{,}00}
% alignmentgraph-isd run-to-run SD (SD of the run-level means) at qwen3.5-9b
\newcommand{\rqOneHarnessRunSDNineB}{0{,}40}
% best baseline (dick-carey-agent) pooled mean Total at qwen3.5-9b
\newcommand{\rqOneBestBaselineTotalNineB}{77{,}00}
% name of the best baseline at qwen3.5-9b
\newcommand{\rqOneBestBaselineNameNineB}{Dick-Carey-Agent}
% mean paired delta alignmentgraph-isd - dick-carey-agent at qwen3.5-9b
\newcommand{\rqOneDeltaBestNineB}{4{,}00}
% bootstrap 95% CI lower bound of that delta
\newcommand{\rqOneDeltaBestCILoNineB}{3{,}00}
% bootstrap 95% CI upper bound of that delta
\newcommand{\rqOneDeltaBestCIHiNineB}{5{,}00}
% max Holm-adjusted p over the 6 comparisons at qwen3.5-9b
\newcommand{\rqOnePHolmMaxNineB}{6{,}0\times 10^{-9}}
% max Holm-adjusted p across ALL sizes x baselines (complete-case)
\newcommand{\rqOnePHolmMaxAll}{6{,}0\times 10^{-9}}
% alignmentgraph-isd Total at largest size minus smallest size (qwen3.5-9b - qwen3.5-0.8b); positive = larger is better
\newcommand{\rqOneHarnessDropSmall}{35{,}00}
% DiD of delta(proposed - baseline): largest minus smallest size
\newcommand{\rqOneDiDVsBaseline}{-2{,}00}
% bootstrap 95% CI lower bound of that DiD
\newcommand{\rqOneDiDVsBaselineCILo}{-3{,}00}
% bootstrap 95% CI upper bound of that DiD
\newcommand{\rqOneDiDVsBaselineCIHi}{-1{,}00}
% Wilcoxon p of the per-scenario DiD values vs 0
\newcommand{\rqOneDiDVsBaselineP}{0{,}010}
% OLS slope of delta(proposed - baseline) per size step (4 sizes)
\newcommand{\rqOneTrendSlopeVsBaseline}{-0{,}66}
% sentinel: nonempty means the generated numbers are placeholders
\newcommand{\ladderDataIsDemo}{TBD-DEMO}


% Shared assets (references.bib, figures/) live in the parent acm/ directory.
\graphicspath{{../figures/}{figures/}}

% Bản dịch tiếng Việt của main.tex --- xem CLAUDE.md và main.tex (bản gốc tiếng Anh)
% cho các TODO trước khi nộp (venue, ISBN/DOI, CCS codes).
% TODO camera-ready: đổi thành printacmref=true và điền venue/ISBN/DOI thật.
\settopmatter{printacmref=false}
\renewcommand\footnotetextcopyrightpermission[1]{}
\acmConference[Preprint]{Preprint --- venue TBD}{2026}{}
\acmYear{2026}
\acmISBN{}
\acmDOI{}
\setcopyright{none}

\title{AlignmentGraph-ISD: A Graph-Guided Multi-Agent Harness for Instructional Systems Design with Small Language Model}

% TODO: xác nhận email và tên đơn vị công tác chính xác trước khi nộp.
\author{Phuong Viet Bui}
\affiliation{%
  \department{Trường Công nghệ Thông tin và Truyền thông}
  \institution{Đại học Bách khoa Hà Nội}
  \city{Hà Nội}
  \country{Việt Nam}
}
\email{phuong.bv242011m@sis.hust.edu.vn}

\author{Hai Van Pham}
\authornote{Giảng viên hướng dẫn chính.}
\affiliation{%
  \department{Trường Công nghệ Thông tin và Truyền thông}
  \institution{Đại học Bách khoa Hà Nội}
  \city{Hà Nội}
  \country{Việt Nam}
}
\email{haipv@soict.hust.edu.vn}

%\author{Long Kim Cu}
%\authornote{Giảng viên đồng hướng dẫn.}
%\affiliation{%
%  \institution{Center for Information Technology, Ministry of Science and Technology}
%  \city{Hanoi}
%  \country{Vietnam}
%}

\renewcommand{\shortauthors}{Bui và Pham}

% Cho phép TeX co giãn thêm ở bước cuối để tránh overfull hbox do các token dài
% không ngắt được (vd AlignmentGraph-ISD, verification/repair) trong chế độ tiếng Việt.
\setlength{\emergencystretch}{2em}

\begin{document}

\begin{abstract}
Các agent dựa trên LLM (large language model) có thể sinh trọn vẹn một chương
trình giảng dạy (instructional program), nhưng chất lượng của chúng phụ thuộc vào
\emph{constructive alignment}---sự tương ứng giữa learning objective, activity và
assessment---vốn là thuộc tính sụp đổ đầu tiên khi các tổ chức chuyển từ model lớn
sang small language model (SLM) chi phí thấp để sinh nội dung đào tạo nội bộ. Suy
giảm alignment là một khiếm khuyết tích lũy và vô hình: chương trình vẫn "nhìn ổn"
theo từng phần và chỉ lộ lỗi khi triển khai. Các hướng LLM-agent hiện có cho
instructional systems design (ISD) chỉ \emph{khuyến khích} alignment qua prompt;
không hướng nào biểu diễn nó thành trạng thái mà hệ thống có thể kiểm tra và sửa,
nên chất lượng bị phó mặc cho năng lực thô của model. Chúng tôi đề xuất
\textbf{AlignmentGraph-ISD}, một graph-guided multi-agent harness hiện thực hoá
constructive alignment thành một \emph{alignment graph} tường minh và kiểm tra được,
rồi tổ chức việc sinh xoay quanh nó: các agent chuyên biệt xây dựng graph theo một workflow cố định với
phạm vi ghi bị giới hạn, nên alignment hình thành \emph{by construction}, và mỗi
agent chạy pedagogical verification cùng repair như một công cụ ngay trong lượt của
mình. Chúng tôi đánh giá harness trên ISD-Agent-Bench, so với sáu agent tiêu biểu,
dưới cả một model lớn lẫn một SLM 9B, và thực hiện một ablation để tách riêng đóng
góp của từng cơ chế. Trên benchmark, AlignmentGraph-ISD đạt điểm tổng cao nhất
trong bảy agent; ablation cho thấy multi-agent decomposition trên graph là cơ chế
đóng góp chính ($+3{,}9$ điểm tổng, $p < 10^{-12}$), và---quan trọng---đóng
góp này gần như gấp đôi khi model nhỏ đi (từ $+3{,}9$ điểm tổng trên LLM lên $+7{,}9$
trên SLM, $p < 10^{-12}$), trong khi điểm của chính harness gần như bất biến theo
scale còn các baseline suy giảm. Pedagogical verification/repair không làm thay đổi
điểm do judge chấm trên cả hai model dù giảm được số alignment defect đo được. Các
kết quả này cho thấy: biểu diễn alignment thành trạng thái graph hạng nhất, và phân
rã việc sinh xoay quanh nó, chính là điều cho phép một agent harness duy trì chất
lượng thiết kế giảng dạy độc lập với model bên dưới.
\end{abstract}


\begin{CCSXML}
<ccs2012>
   <concept>
       <concept_id>10010405.10010489</concept_id>
       <concept_desc>Applied computing~Education</concept_desc>
       <concept_significance>500</concept_significance>
   </concept>
   <concept>
       <concept_id>10010147.10010178.10010219.10010220</concept_id>
       <concept_desc>Computing methodologies~Multi-agent systems</concept_desc>
       <concept_significance>500</concept_significance>
   </concept>
   <concept>
       <concept_id>10010147.10010178.10010179</concept_id>
       <concept_desc>Computing methodologies~Natural language processing</concept_desc>
       <concept_significance>300</concept_significance>
   </concept>
</ccs2012>
\end{CCSXML}

\ccsdesc[500]{Applied computing~Education}
\ccsdesc[500]{Computing methodologies~Multi-agent systems}
\ccsdesc[300]{Computing methodologies~Natural language processing}

\keywords{instructional systems design, instructional design, LLM agent, agent harness, constructive alignment, multi-agent system, small language model}

\maketitle

\section{Giới thiệu}
\label{sec:introduction}

Trong những năm gần đây, các large language model (LLM) thương mại đã tiến từ chỗ hỗ
trợ soạn thảo từng phần sang đảm đương gần trọn quy trình instructional systems
design (ISD): các hệ thống dựa trên GPT-4 và những model cùng lớp có thể sinh được
lesson plan, bộ câu hỏi đánh giá, thậm chí trọn gói học liệu của một môn học từ một
yêu cầu thiết kế duy
nhất~\cite{yao2025instructionalagents,zhang2025eduplanner,agentlesson2026,zeng2026coma,hu2024teachingplan}.
Chất lượng ở cấp thành phần đã được kiểm chứng qua nhiều nghiên cứu so sánh:
learning objective tương đương bản do giảng viên soạn~\cite{doyle2025objectives},
câu hỏi trắc nghiệm ngang bản người viết~\cite{doughty2024mcq}, giáo án không phân
biệt được với bản của đồng nghiệp khi chấm
blind~\cite{cooper2026indistinguishable}, và hiệu quả lớp học không thua
kém~\cite{karaman2024chatgptlesson}; với vòng lặp mô phỏng--tinh chỉnh, chất lượng
đạt mức của giáo viên giàu kinh nghiệm~\cite{hu2025teachingsim}. Các công cụ này
cũng đã đi vào thực hành của instructional
designer~\cite{yang2025idpractice,lessonplanner2024}.

Tuy nhiên, việc triển khai các hệ thống này tại cơ sở giáo dục gặp trở ngại vì
chúng phụ thuộc vào API thương mại. Một quy trình ISD hoàn chỉnh cần nhiều lời gọi model nên
chi phí tích luỹ lớn; khảo sát trong giáo dục đại học ghi nhận ràng buộc tài chính
đang làm chậm việc ứng dụng AI~\cite{robert2025educause}, còn dữ liệu người học
phải đi qua hạ tầng bên thứ ba trong khi vấn đề riêng tư được đánh giá là chưa xử
lý thoả đáng~\cite{yan2024ethical}. Nhiều đại học đã dựng gateway riêng để kiểm
soát ngân sách và dữ liệu~\cite{mithun2025aiverde}; tự vận hành model open-weight
cũng rẻ hơn API khi khối lượng sử dụng đủ lớn~\cite{pan2025onpremise}. Như vậy,
chi phí và an toàn dữ liệu được ghi nhận là rào cản triển khai tại các cơ sở giáo
dục, và hướng small language model (SLM) open-weight chạy on-premise đang được
quan tâm: hạ tầng LLM mở cho đại học đào tạo từ xa~\cite{zesch2024flexi}, đánh giá
model tự host cho trường học~\cite{chan2026bigger}, khung biên soạn nội dung chạy
cục bộ vì riêng tư và chi phí~\cite{reza2025smallmodels}.

Khi dùng SLM, model nền yếu hơn, nên chất lượng đầu ra phụ thuộc nhiều hơn vào
harness bao quanh, tức là cách công việc được phân rã, điều phối và kiểm soát chất
lượng. Model nhỏ yếu khi dùng tool nhưng khung đa tác tử có phân rã thu hẹp khoảng
cách với model lớn~\cite{shen2024multillmagent}; scaffolding dạng agent giúp model
0.5B--3B đạt gần chất lượng của nhóm lớn hơn~\cite{kang2025agentdistill}; SLM cũng
được lập luận là đủ cho phần lớn agentic
workload~\cite{belcak2025slm,sharma2025slmagents}. Tuy nhiên, độ nhạy với harness
không đơn điệu giữa các mức kích thước model~\cite{cho2026harness}, nên quan hệ
giữa chất lượng, harness và kích thước model cần được đo lường thay vì giả định.

Từ bối cảnh trên, khoảng trống nghiên cứu trung tâm hiện ra: các hệ tác tử ISD và
lesson-plan hiện có đều được xây dựng và đánh giá trên model thương mại mạnh, còn
các công trình gần nhất mỗi công trình chỉ chạm một phần của bài
toán~\cite{jaldi2026slmassessment,zheng2025lessonplanlm,reza2025smallmodels,lee2026educlaw,wei2026slmeducation}.
Theo khảo sát của chúng tôi, chưa có công trình nào cho thấy một hệ tác tử duy
trì được chất lượng ISD end-to-end trên model nhỏ mà không cần fine-tune, cũng
như chưa có đánh giá hệ thống nào về chất lượng đó trên một dải kích thước model.

Đồng thời, khoảng trống đó cũng làm nảy sinh một vấn đề khác liên quan đến phương pháp đo lường chất lượng,
bởi hai công cụ thông dụng đều có giới hạn. Rubric chấm bằng LLM-judge --- cách của các
benchmark ISD và giáo dục hiện
nay~\cite{isdagentbench,edubench2025,tutorbench2025,mathtutorbench2025} --- mang
các bias đã được ghi nhận có hệ
thống~\cite{zheng2023judging,panickssery2024llm,norman2026reliability}, và ở đây
họ model chấm còn trùng với họ model sinh. Đánh giá bằng chuyên gia con người thì
đáng tin hơn nhưng không scale cho một ladder gồm nhiều kích thước, nhiều run và
nhiều hệ. Quan trọng hơn cả, constructive alignment đặt chất lượng của một thiết
kế ở quan hệ giữa mục tiêu, hoạt động và đánh giá~\cite{biggs1996enhancing}: một
thiết kế có thể đọc trôi chảy, nhìn qua có vẻ tốt, nhưng vẫn sai lệch nghiêm trọng về
alignment. Không cách đo nào ở trên làm tường minh được các tín hiệu alignment
ấy. Các phương pháp đo alignment tự động hiện có cũng chỉ đo sự ăn khớp
với một bộ chuẩn hoặc tham chiếu soạn
sẵn~\cite{porter2002measuring,webb1997criteria,xu2025itemalignment,silva2026lessonbench},
chưa đo alignment nội tại của một thiết kế hoàn toàn do máy sinh.

Mục tiêu của nghiên cứu này là duy trì chất lượng thiết kế ISD khi thu nhỏ model
bằng một harness đa tác tử có trạng thái alignment tường minh
(\Cref{sec:method}), đồng thời làm tường minh các tín hiệu alignment của thiết kế
được sinh bằng một giao thức đo không dùng LLM (\Cref{sec:protocol}). Mục tiêu đó
được cụ thể hoá thành hai câu hỏi nghiên cứu:
\begin{itemize}
  \item \textbf{RQ1 (duy trì chất lượng theo thang LLM-as-judge).} Khi model nền
  nhỏ dần trên ladder Qwen 0.8B--9B, AlignmentGraph-ISD có duy trì được chất lượng
  thiết kế đo bằng LLM-as-judge của ISD-Agent-Bench so với sáu benchmark agent
  không: hệ đứng ở đâu tại từng kích thước, và khoảng cách có nới rộng khi model
  nhỏ đi không?
  \item \textbf{RQ2 (duy trì chất lượng theo thang tín hiệu alignment tường minh).} Khi được làm tường
  minh bằng giao thức đo không dùng LLM, các tín hiệu alignment nội tại của thiết
  kế do máy sinh thay đổi thế nào theo kích thước model, và khoảng cách so sánh có
  nới rộng khi model nhỏ đi không?
\end{itemize}

Bài báo có hai đóng góp chính:
\begin{enumerate}
  \item \textbf{AlignmentGraph-ISD}: một harness đa tác tử tổ chức việc sinh
  chương trình ADDIE quanh một alignment graph, tức trạng thái alignment tường
  minh, có thể kiểm tra và audit. Cuối mỗi lần chạy, toàn bộ graph được xuất ra để
  instructional designer xem xét từng liên kết mục tiêu--hoạt động--đánh giá và
  can thiệp khi cần, theo nguyên tắc teacher-in-the-loop.
  \item \textbf{Giao thức đánh giá constructive alignment tự động}, mở rộng
  ISD-Agent-Bench: không dùng LLM trong đo lường, mỗi điểm số là một hàm cố định
  của văn bản, kết quả là panel bảy tín hiệu ngang hàng thuộc hai họ công cụ đo
  độc lập về cơ chế. Giao thức operationalize các tiêu chí đã được thiết lập thay
  vì đề xuất construct
  mới~\cite{biggs1996enhancing,porter2002measuring,webb1997criteria}; phạm vi
  diễn giải giới hạn ở các tín hiệu đo được.
\end{enumerate}

Phần còn lại của bài báo được tổ chức như sau: \Cref{sec:related-work} trình bày
các công trình liên quan; \Cref{sec:method} mô tả AlignmentGraph-ISD và giao thức
đánh giá (\Cref{sec:protocol}); \Cref{sec:experimental-methodology} mô tả thiết kế
thực nghiệm; \Cref{sec:results} báo cáo kết quả; \Cref{sec:discussion} thảo luận
kết quả và các nguy cơ đối với tính hợp lệ; \Cref{sec:conclusion} kết luận bài
báo.

\section{Công trình liên quan}
\label{sec:related-work}

\TBD{}

% ==== BẢN NHÁP CŨ (giữ tham khảo, sẽ viết lại) ====
% \section{Công trình liên quan}
% \label{sec:related-work}
%
% \subsection{Constructive alignment và các khung đo alignment}
%
% Constructive alignment định nghĩa chất lượng giảng dạy qua quan hệ giữa mục tiêu
% học tập, hoạt động dạy--học và assessment, thay vì qua từng artifact riêng
% lẻ~\cite{biggs1996enhancing,biggs2011teaching}. Quality Matters đưa các quan hệ này
% vào quy trình review: objective phải đo được, còn assessment, học liệu và hoạt động
% phải hỗ trợ objective~\cite{qm2018rubric}. Về phía đo lường định lượng, chỉ số
% Porter đo mức tương hợp giữa hai phân bố nội dung--cognitive demand, các tiêu chí
% của Webb đánh giá những mặt như tính nhất quán của cognitive demand, và Fulmer khảo
% sát cách diễn giải các chỉ số
% này~\cite{porter2002measuring,webb1997criteria,fulmer2011estimating}. Bloom
% taxonomy bản sửa đổi cung cấp bộ từ vựng chung về cognitive
% demand~\cite{anderson2001taxonomy}. Điểm chung: các khung này đo trên curriculum,
% bộ chuẩn và assessment do con người soạn sẵn từ trước.
%
% \subsection{Hệ tác tử ISD và benchmark đánh giá}
%
% ISD-Agent-Bench tập hợp các agent xây theo những lý thuyết thiết kế khác nhau, từ
% ADDIE, thiết kế hệ thống, rapid prototyping tới vòng lặp
% reasoning--acting~\cite{isdagentbench,yao2023react}. Nhiều hệ khác hỗ trợ soạn
% giáo án theo định hướng sư phạm hoặc sinh lesson plan bằng cơ chế đa tác
% tử~\cite{lessonplanner2024,hu2024teachingplan,zhang2025eduplanner,agentlesson2026,zeng2026coma}.
% Về đánh giá, ISD-Agent-Bench và các benchmark giáo dục rộng hơn đều chấm bằng
% rubric LLM-judge~\cite{edubench2025,tutorbench2025,mathtutorbench2025}, trong khi
% bias của LLM-judge đã được ghi nhận có hệ
% thống~\cite{zheng2023judging,panickssery2024llm,norman2026reliability}. Xuyên suốt
% nhóm công trình này, chưa hệ nào đặt duy trì chất lượng khi thu nhỏ model làm mục
% tiêu trung tâm, và chưa benchmark nào instrument riêng các quan hệ giữa objective,
% hoạt động và assessment do máy sinh.
%
% \subsection{Các tiền lệ gần nhất}
%
% Về phía triển khai, các công trình gần nhất mỗi công trình giải quyết một phần bài
% toán. Small Models, Big Support chạy model self-hosted để tạo nội dung và
% assessment cho giáo viên, nhưng không có quy trình ISD agentic đầu
% cuối~\cite{reza2025smallmodels}. Jaldi và cộng sự so sánh SLM với LLM cho khâu sinh
% câu hỏi đánh giá với động cơ triển khai bảo vệ riêng tư, nhưng dừng ở một thành
% phần đơn lẻ~\cite{jaldi2026slmassessment}. LessonPlanLM đạt chất lượng lesson plan
% trên model 7B bằng fine-tune, thay vì kiểm tra một harness qua nhiều kích
% thước~\cite{zheng2025lessonplanlm}. Instructional Agents sinh trọn bộ học liệu và
% báo cáo chi phí, nhưng vẫn chạy trên API thương
% mại~\cite{yao2025instructionalagents}. EduClaw-Bench khảo sát tương tác giữa thiết
% kế agent và tier model, nhưng cho bài toán tutoring và vẫn chấm bằng
% LLM-judge~\cite{lee2026educlaw}. Như vậy, authoring tại chỗ, chuyên biệt hoá model
% nhỏ và sinh học liệu đa tác tử đã được nghiên cứu riêng lẻ~\cite{wei2026slmeducation};
% chúng chưa được kết hợp thành một phép đo có kiểm soát cho một ISD harness duy
% nhất trên dải model open-weight.
%
% Về phía đo lường, các thành phần cần thiết đều đã tồn tại. Semantic traceability và
% embedding đa ngôn ngữ khôi phục quan hệ giữa các artifact văn
% bản~\cite{clelandhuang2018traceability,feng2022labse}. Embedding và LM được dùng để
% phân loại assessment item vào bộ chuẩn nội dung có
% sẵn~\cite{xu2025itemalignment}. LessonBench-V1 so phân bố Bloom của bài giảng máy
% sinh với lesson plan tham chiếu đã qua duyệt của con
% người~\cite{silva2026lessonbench}. Với cognitive demand, Li và cộng sự công bố
% corpus learning objective do chuyên gia gán nhãn; nghiên cứu này dùng corpus đó để
% fine-tune Bloom classifier~\cite{li2022bloomclassification}, còn BloomNet là một
% classifier liên quan chứ không phải nguồn corpus~\cite{waheed2021bloomnet}. Các
% phương pháp trên không cần generative judge, nhưng đều đo sự ăn khớp với một bộ
% chuẩn hoặc tham chiếu có sẵn; chưa phương pháp nào đánh giá đồng thời các artifact
% do chính hệ bị đo sinh ra.
%
% \subsection{Tổng hợp}
%
% \Cref{tab:related} định vị các nhóm công trình theo tiêu chí; bảng mô tả phạm vi,
% không xếp hạng chất lượng. Hai tổ hợp còn bỏ ngỏ hiện rõ: đánh giá có kiểm soát
% theo scale cho một ISD harness chạy tại chỗ, và thước đo alignment nội tại tích
% hợp vào benchmark ISD --- nghiên cứu này xử lý tổ hợp thứ nhất bằng ladder Qwen
% 0.8B--9B và tổ hợp thứ hai bằng giao thức panel bảy tín hiệu.
%
% \begin{table}[t]
%   \caption{Định vị công trình liên quan theo tiêu chí. ISD = instructional design
%   agentic đầu cuối; Local = SLM local/open-weight là điều kiện thiết kế; Scale = đánh
%   giá có kiểm soát qua các kích thước model; Intr. = alignment nội tại giữa các
%   design artifact đều do máy sinh; Bench = tích hợp vào ISD benchmark.
%   \ding{51} = đã xét; \ding{55} = chưa xét; (\ding{51}) = đã xét một phần.}
%   \label{tab:related}
%   \centering
%   \small
%   \resizebox{\columnwidth}{!}{%
%   \begin{tabular}{@{}lccccc@{}}
%     \toprule
%     Công trình hoặc nhóm & ISD & Local & Scale & Intr. & Bench \\
%     \midrule
%     ISD-Agent-Bench và các agent~\cite{isdagentbench} & \ding{51} & \ding{55} & \ding{55} & \ding{55} & \ding{55} \\
%     Các hệ sinh lesson plan bằng LLM~\cite{lessonplanner2024,hu2024teachingplan,zhang2025eduplanner,agentlesson2026,zeng2026coma} & (\ding{51}) & \ding{55} & \ding{55} & \ding{55} & \ding{55} \\
%     Small Models, Big Support~\cite{reza2025smallmodels} & \ding{55} & \ding{51} & \ding{55} & \ding{55} & \ding{55} \\
%     SLM cho sinh câu hỏi đánh giá~\cite{jaldi2026slmassessment} & \ding{55} & \ding{51} & (\ding{51}) & \ding{55} & \ding{55} \\
%     LessonPlanLM~\cite{zheng2025lessonplanlm} & \ding{55} & (\ding{51}) & \ding{55} & \ding{55} & \ding{55} \\
%     Instructional Agents~\cite{yao2025instructionalagents} & (\ding{51}) & \ding{55} & \ding{55} & \ding{55} & \ding{55} \\
%     EduClaw-Bench~\cite{lee2026educlaw} & \ding{55} & \ding{55} & (\ding{51}) & \ding{55} & \ding{55} \\
%     Đo item--standards bằng LM/embedding~\cite{xu2025itemalignment} & \ding{55} & \ding{55} & \ding{55} & \ding{55} & \ding{55} \\
%     LessonBench-V1~\cite{silva2026lessonbench} & (\ding{51}) & \ding{55} & \ding{55} & \ding{55} & (\ding{51}) \\
%     Bloom classifier dựa trên BERT~\cite{li2022bloomclassification,waheed2021bloomnet} & \ding{55} & \ding{55} & \ding{55} & \ding{55} & \ding{55} \\
%     \textbf{AlignmentGraph-ISD} (nghiên cứu này) & \ding{51} & \ding{51} & \ding{51} & \ding{51} & \ding{51} \\
%     \bottomrule
%   \end{tabular}}
% \end{table}

\section{Mô hình hệ thống và phát biểu bài toán}
\label{sec:problem-formulation}

Khoảng trống nghiên cứu đặt ra câu hỏi rộng về khả năng duy trì chất lượng ISD khi mô
hình nền bị thu nhỏ. Để chuyển câu hỏi này thành một đối tượng mà harness có thể thao
tác, chúng tôi dùng constructive alignment để vận hành hoá sự nhất quán giữa các thành
phần của thiết kế giảng dạy.
Lựa chọn này không đồng nhất toàn bộ chất lượng ISD với alignment; nó cung cấp một thuộc
tính quan hệ trung tâm để hình thức hoá cơ chế và kiểm tra giả thuyết thiết kế.

\subsection{Mô hình hệ thống}
\label{sec:system-model}

\Cref{fig:architecture} mô tả ranh giới trách nhiệm giữa mô hình và harness. Mô hình
ngôn ngữ sinh nội dung; harness duy trì trạng thái, điều phối quy trình và thực hiện
kiểm soát chất lượng. Mọi thành phần của harness cùng đọc hoặc cập nhật một alignment
graph có phiên bản, đóng vai trò trạng thái thiết kế trung tâm.

\begin{figure*}[t]
  \centering
  \includegraphics[width=\textwidth]{figures/architecture}
  \Description{Sơ đồ kiến trúc AlignmentGraph-ISD. Các tác tử chuyên biệt chạy theo
  thứ tự cố định và chỉ ghi những kiểu node, edge thuộc phạm vi của mình vào alignment
  graph. Sau mỗi lần ghi, tác tử kiểm chứng và sửa phần trạng thái do nó sở hữu. Bộ tổng
  hợp cuối cùng kết xuất chương trình ADDIE cùng artifact đồ thị phục vụ audit.}
  \caption{Kiến trúc AlignmentGraph-ISD. Các tác tử chuyên biệt lần lượt xây dựng
  alignment graph dưới phạm vi ghi giới hạn, rồi kiểm chứng và sửa lỗi ngay trong lượt
  của mình. Bộ tổng hợp cuối kết xuất chương trình ADDIE cùng artifact đồ thị auditable
  (\cref{sec:aggregation}). Mô hình nền có thể được thay thế mà không làm thay đổi các
  cơ chế của harness.}
  \label{fig:architecture}
\end{figure*}

\subsection{Alignment Graph}
\label{sec:alignment-graph}

Trong đầu ra dạng văn bản, quan hệ giữa các thành phần thiết kế chỉ tồn tại ngầm: người
đọc có thể nhận ra một assessment đang đo mục tiêu nào, nhưng hệ thống không có đối
tượng để kiểm tra quan hệ đó. AlignmentGraph-ISD giải quyết vấn đề này bằng cách hiện
thực hoá constructive alignment~\cite{biggs1996enhancing} thành cấu trúc dữ liệu tường
minh. Tại bước $t$, trạng thái thiết kế là đồ thị $G_t=(V_t,E_t)$. Tập đỉnh bao phủ các
thành phần thuộc năm pha ADDIE:
\begin{multline*}
V_t = V_{\mathrm{needs}} \cup V_{\mathrm{objective}} \cup V_{\mathrm{assessment}} \cup V_{\mathrm{activity}} \\
 \cup V_{\mathrm{material}} \cup V_{\mathrm{implementation}} \cup V_{\mathrm{evaluation}}.
\end{multline*}
Tập cạnh chỉ giữ ba quan hệ nhiều--nhiều trực tiếp phục vụ alignment:
\begin{align*}
E_t &= E_{\mathrm{measures}} \cup E_{\mathrm{supports}} \cup E_{\mathrm{evaluates}},\\
E_{\mathrm{measures}} &\subseteq V_{\mathrm{assessment}} \times V_{\mathrm{objective}},\\
E_{\mathrm{supports}} &\subseteq V_{\mathrm{activity}} \times V_{\mathrm{objective}},\\
E_{\mathrm{evaluates}} &\subseteq V_{\mathrm{evaluation}} \times V_{\mathrm{objective}}.
\end{align*}
Cạnh \textsl{measures} nối một assessment với mục tiêu mà nó đo; cạnh
\textsl{supports} nối một hoạt động với mục tiêu mà nó luyện tập; cạnh
\textsl{evaluates} nối một evaluation item với mục tiêu mà nó thẩm định. Quan hệ thứ ba
mở rộng bộ đôi activity--assessment quen thuộc để phù hợp với pha Evaluation của ADDIE:
việc đánh giá chương trình cũng phải quay lại các mục tiêu đã đặt ra. Các node
\texttt{needs}, \texttt{material} và \texttt{implementation} vẫn cần cho đầu ra ADDIE
hoàn chỉnh, nhưng không tham gia trực tiếp vào vị từ alignment; chúng được kiểm bằng các
ràng buộc khác, chẳng hạn tính khả thi của kế hoạch triển khai.

Mỗi mục tiêu học tập được biểu diễn bằng một node riêng
(\texttt{objective:LO-*}), nhờ đó hệ thống có thể kiểm tra alignment ở đúng đơn vị cần
được bảo đảm. Với mọi $o\in V_{\mathrm{objective}}$, điều kiện cấu trúc là:
\begin{multline}
\label{eq:aligned}
\mathrm{aligned}(o) \iff \big(\exists a\colon (a,o) \in E_{\mathrm{measures}}\big) \\
 \wedge \big(\exists x\colon (x,o) \in E_{\mathrm{supports}}\big) \wedge \big(\exists e\colon (e,o) \in E_{\mathrm{evaluates}}\big).
\end{multline}
\Cref{eq:aligned} mới xác nhận \emph{độ phủ cấu trúc}: mỗi mục tiêu có ít nhất một
liên kết thuộc từng loại. Constructive alignment còn đòi hỏi các liên kết đó hợp lệ về
nội dung và cấp độ nhận thức. Ví dụ, một assessment chỉ nhận biết khái niệm không thể
được xem là phép đo đầy đủ cho một mục tiêu yêu cầu phân tích. Verifier chịu trách nhiệm
kiểm tra lớp hợp lệ này (\cref{sec:qc}). Sau khi kiểm, cạnh được gắn cờ
\texttt{verified}; khi một node bị sửa, các node phụ thuộc có thể được đánh dấu
\texttt{stale} để buộc kiểm lại. Các trường này lưu dấu vết kiểm soát chất lượng trong
trạng thái chung.

Vị từ trong \cref{eq:aligned} đồng thời là tiêu chí thành công của bài toán, nên nó
cần được đo bằng một phép đo không phụ thuộc vào cơ chế nội bộ của harness. Thước đo
alignment trong \cref{sec:alignment-metric} lượng hoá trực tiếp cả hai lớp của vị từ
này---độ phủ cấu trúc và tính hợp lệ về nội dung, cấp độ nhận thức---trên đầu ra cuối
của \emph{mọi} hệ được so sánh, chỉ từ văn bản kết xuất. Nhờ đó, đại lượng được dùng để
định nghĩa thành công cũng chính là đại lượng được đo trong thực nghiệm (RQ2), thay vì
chỉ được xấp xỉ gián tiếp qua rubric chất lượng tổng quát.

\subsection{Ký hiệu và giả định}
\label{sec:assumptions}

\Cref{tab:notation} tóm tắt các ký hiệu. Chúng tôi giả định rằng: (i) design brief $b$
đã được cung cấp, còn quá trình hình thành brief nằm ngoài phạm vi; (ii) mô hình nền
không lưu trạng thái giữa các lượt và có thể được thay thế; (iii) đầu ra phải tuân theo
schema ADDIE của benchmark. Vì vậy, ADDIE là hợp đồng đầu ra, không phải yêu cầu duy nhất
đối với thứ tự thực thi bên trong harness.

\begin{table}[t]
  \caption{Ký hiệu.}
  \label{tab:notation}
  \small
  \begin{tabular}{@{}ll@{}}
    \toprule
    Ký hiệu & Ý nghĩa \\
    \midrule
    $b$ & design brief (input của harness) \\
    $G_t=(V_t,E_t)$ & alignment graph tại bước $t$ \\
    $o \in V_{\mathrm{objective}}$ & một learning objective (element node) \\
    $\mathrm{aligned}(o)$ & vị từ constructive-alignment, \cref{eq:aligned} \\
    $A_i$ & specialist agent cho concern ISD $i$ \\
    $V^{(i)},E^{(i)}$ & kiểu node/edge agent $i$ được ghi \\
    $\mathcal{V}$ & pedagogical verifier; $\Omega=\mathcal{V}(G)$ tập vi phạm \\
    $\rho,\gamma,\mu$ & hàm lấy context từ brief, graph, và bộ nhớ domain \\
    $M$ & bộ nhớ domain ổn định (ADDIE, Bloom, rubric template) \\
    $B$ & ngân sách lặp repair mỗi agent (\cref{sec:qc}) \\
    \bottomrule
  \end{tabular}
\end{table}

\subsection{Phát biểu bài toán}
\label{sec:problem-statement}

Với design brief $b$, mục tiêu là sinh một chương trình giảng dạy hoàn chỉnh, gồm đồ thị
$G$ và các thành phần ADDIE được kết xuất, sao cho:
\begin{enumerate}
  \item[(i)] mọi mục tiêu đều đạt độ phủ cấu trúc theo \cref{eq:aligned}, đồng thời
  các cạnh liên quan hợp lệ về nội dung và cấp độ nhận thức; và
  \item[(ii)] cơ chế duy trì chất lượng nằm trong harness và trạng thái alignment là
  tường minh, kiểm tra được bởi cả máy lẫn con người, với mục tiêu hạn chế mức suy giảm
  chất lượng khi mô hình nền được thu nhỏ. Mức độ đạt được mục tiêu này được kiểm tra
  bằng thực nghiệm trong \cref{sec:results}.
\end{enumerate}

\subsection{Thách thức}
\label{sec:challenges}

Bài toán khó vì ba lý do. Thứ nhất, alignment liên kết nhiều thành phần: không tác tử
nào chỉ nhìn phần đầu ra của mình có thể xác nhận toàn bộ điều kiện. Thứ hai, vi phạm
alignment thường không xuất hiện như một lỗi cục bộ; từng đoạn vẫn hợp lý dù các đoạn
không còn nhất quán với nhau. Thứ ba, nếu cơ chế kiểm soát chỉ nằm trong prompt, chất
lượng tiếp tục phụ thuộc vào khả năng tự tuân thủ của mô hình. Do đó, harness cần một
trạng thái chung có thể kiểm tra và một quy trình cập nhật tường minh.

\section{Harness AlignmentGraph-ISD}
\label{sec:method}

\subsection{Tổng quan}
\label{sec:overview}

Bài toán ở \cref{sec:problem-statement} đòi hỏi alignment phải được sinh ra và đảm
bảo độc lập với model. Quan sát dẫn dắt thiết kế của chúng tôi là: mọi thứ điều này
cần---trạng thái bền vững, kiểm soát chất lượng, và điều phối chuyên biệt---là trách
nhiệm của môi trường thực thi, không phải năng lực của model. Do đó AlignmentGraph-ISD
tập trung các trách nhiệm này vào một harness, để model chỉ sinh nội dung bên trong
nó. Một agent hoàn chỉnh là sự kết hợp của cả hai (\emph{agent = harness + model}).
Cụ thể, harness (i)~phân rã việc sinh cho các agent chuyên biệt xây dựng alignment
graph theo một thứ tự thiết kế giảng dạy cố định với phạm vi ghi giới hạn
(\cref{sec:decomposition}), để alignment hình thành by construction, và (ii)~trao
cho mỗi agent pedagogical verification cùng repair như các công cụ nó chạy ngay
trong lượt của mình (\cref{sec:qc}), để vi phạm được bắt và sửa trước khi việc sinh
tiếp tục. Một bước aggregation cuối (\cref{sec:aggregation}) phát ra chương trình và
ghi lại một residual audit.

\subsection{Multi-Agent Decomposition trên Graph}
\label{sec:decomposition}

Mỗi concern ISD do một specialist agent xử lý---needs context, objective design,
assessment design, activity và material, implementation plan, và evaluation---chạy
theo thứ tự ADDIE cố định. Mỗi agent $A_i$ là một toán tử cập-nhật-graph có phạm vi
ghi bị giới hạn bởi chuyên môn của nó:
\[
A_i : (G_t, c) \mapsto G_{t+1}, \qquad \Delta(G_t, G_{t+1}) \subseteq V^{(i)} \cup E^{(i)},
\]
trong đó $\Delta(G_t, G_{t+1})$ là các node và edge được thêm giữa hai graph, còn
$V^{(i)}, E^{(i)}$ là các kiểu node và edge mà agent $i$ được ghi. Ví dụ, assessment
agent chỉ ghi node \texttt{assessment} và edge \textsl{measures}: khi nó phát ra một
assessment item nó phải nêu tên objective mà item chấm, và việc nêu tên đó \emph{chính
là} edge \textsl{measures}. Vì objective agent chạy trước và đặt các objective vào
graph chung, mọi agent sau tham chiếu tới objective đã có, và các edge alignment của
\cref{eq:aligned} được tạo ra như một phần bắt buộc của việc sinh. Nhờ vậy alignment
đúng \emph{by construction}, thay vì được kỳ vọng trong văn bản. Context mỗi agent
nhận là $c = \rho(b) \cup \gamma(G_t) \cup \mu(M)$, với $\rho(b)$ là context từ
brief, $\gamma(G_t)$ là graph-context đọc từ các node trước, và $\mu(M)$ là bộ nhớ
domain ổn định (ADDIE, Bloom, rubric template).

\subsection{Verification và Repair như Công cụ Per-Agent}
\label{sec:qc}

Chỉ riêng construction đảm bảo các edge \emph{tồn tại}; nó không đảm bảo mỗi edge
\emph{hợp lệ} hay mọi objective được phủ. Đó là việc của verification và repair.
Quan trọng: đây không phải một giai đoạn tổng thể riêng chạy sau khi sinh xong; chúng
là các công cụ mỗi specialist chạy \emph{ngay trong lượt của mình}, ngay sau khi viết
artifact. Điều này giữ cho trajectory của mỗi agent là một vòng sạch
generate\,$\rightarrow$\,self-check\,$\rightarrow$\,emit và cho phép các kiểm tra quan
hệ hoạt động, vì đến lúc một agent self-check, các artifact trước mà nó phụ thuộc đã
có trong graph chung.

\paragraph{Verifier.} Gọi $\mathcal{V}$ là pedagogical verifier và
$\Omega = \mathcal{V}(G)$ là tập vi phạm nó phát hiện. Nó chạy bảy họ ràng buộc tất
định trên graph: objective coverage (sự tồn tại
\textsl{measures}/\textsl{supports}/\textsl{evaluates} theo từng objective), tính
hợp lệ của Bloom-verb và nhất quán verb--level~\cite{anderson2001taxonomy}, khớp
\emph{cấp độ nhận thức} assessment--objective (một objective bậc cao chỉ đo bằng item
mức nhận diện sẽ bị đánh dấu), nhất quán rubric (tham chiếu objective treo, tiêu chí
chấm thiếu hoặc không quan sát được), phủ luyện-tập-chủ-động cho objective bậc cao,
tính khả thi triển khai so với brief, và tính cụ thể của formative evaluation. Một
kiểm tra---\emph{measurability} của objective---được báo cho audit nhưng không tự
sửa, vì một Bloom verb quan sát được đã làm objective đo được và việc viết lại các
phát biểu vốn đã tốt sẽ làm giảm chất lượng (\cref{sec:discussion}).

\paragraph{Repair có chủ sở hữu.} Mỗi loại vi phạm thuộc về specialist mà artifact
của nó có thể sửa. Sau khi agent $A_i$ sinh, nó verify và repair chỉ các loại nó sở
hữu, tối đa một budget $B$:
\[
\text{lặp } \big[\, \Omega_i \leftarrow \{\omega \in \mathcal{V}(G) : \mathrm{owner}(\omega)=i\};\ G \leftarrow A_i^{\mathrm{repair}}(G, \Omega_i)\,\big]
\]
cho đến khi $\Omega_i = \emptyset$ hoặc hết budget. Một repair kết hợp một revision
nhắm-đích bằng LLM lên artifact bị đánh dấu với một fallback tất định đảm bảo hội tụ
ngay cả khi output của model không dùng được. Vì repair thuộc về chính specialist đã
sinh, việc bỏ decomposition (\cref{sec:experimental-methodology}) tất yếu bỏ luôn
specialist repair: một biến thể single-agent tự sửa toàn cục cho chính nó, nên không
specialist nào bị lén đưa trở lại.

\subsection{Aggregation và Residual Audit}
\label{sec:aggregation}

Một aggregator cuối kết xuất graph thành ADDIE output contract và kiểm nó theo
schema. Nó cũng ghi một \emph{residual audit im lặng}: chạy lại verifier một lần và
lưu mọi vi phạm còn lại làm metric, không sửa. Điều này phơi bày, với mỗi cấu hình,
bao nhiêu alignment defect còn sót---một thước đo điều kiện (i) độc lập với điểm
rubric dùng ở \cref{sec:results}.

\subsection{Thuật toán}
\label{sec:algorithm}

\begin{figure}[t]
\small
\hrule\vspace{2pt}
\textbf{Thuật toán 1:} Sinh với AlignmentGraph-ISD\vspace{2pt}
\hrule\vspace{4pt}
\begin{tabbing}
xx\=xx\=xx\=\kill
\textbf{Input:} brief $b$; specialist $A_1,\dots,A_k$ theo thứ tự ADDIE; budget $B$\\[1pt]
$G \gets \emptyset$\\
\textbf{for} $i = 1$ \textbf{to} $k$ \textbf{do}\\
\>$c \gets \rho(b) \cup \gamma(G) \cup \mu(M)$\\
\>$G \gets A_i(G, c)$ \quad // sinh; tạo node/edge sở hữu\\
\>\textbf{for} $r = 1$ \textbf{to} $B$ \textbf{do} \quad // self-check trong lượt\\
\>\>$\Omega_i \gets \{\omega \in \mathcal{V}(G) : \mathrm{owner}(\omega)=i\}$\\
\>\>\textbf{if} $\Omega_i = \emptyset$ \textbf{then break}\\
\>\>$G \gets A_i^{\mathrm{repair}}(G, \Omega_i)$\\
\textbf{return} $\mathrm{aggregate}(G)$, \ $\mathrm{residual}\gets\mathcal{V}(G)$
\end{tabbing}
\vspace{-4pt}\hrule
\caption{Sinh với self-check per-agent.}
\label{alg:harness}
\end{figure}

\Cref{alg:harness} tóm tắt việc sinh. Workflow là một pipeline cố định; không dùng
dynamic planner trong setting single-pass nghiên cứu ở đây (một biến thể dùng planner
được giữ lại cho phần mở rộng multi-session ở \cref{sec:threats-limitations}).

\subsection{Khác biệt so với các Baseline Gần nhất}
\label{sec:diff}

So với các baseline mạnh nhất, khác biệt mang tính biểu diễn chứ không chỉ quy trình.
ReAct-ADDIE và các agent bám lý thuyết khác truyền chỉ dẫn alignment qua prompt và
không bao giờ vật chất hoá nó; objective chúng sinh và assessment chúng sinh không
bao giờ được liên kết trong một cấu trúc mà hệ thống có thể kiểm tra.
AlignmentGraph-ISD thay vào đó biến mỗi quan hệ alignment thành một edge có kiểu được
tạo lúc sinh và được kiểm bởi một verifier có chủ sở hữu, nên khác biệt về input,
điều khiển và output là khác biệt về \emph{trạng thái nào tồn tại}: một graph mà
harness có thể cưỡng chế \cref{eq:aligned}, so với văn bản mà nó chỉ có thể prompt.

\section{Phương pháp Thực nghiệm}
\label{sec:experimental-methodology}

\subsection{Câu hỏi Nghiên cứu}
\label{sec:rqs}

\begin{itemize}
  \item \textbf{RQ1 (chất lượng).} Trên ISD-Agent-Bench, AlignmentGraph-ISD so với
  các ISD agent tiêu biểu về chất lượng artifact như thế nào, dưới cả một model lớn
  lẫn một small language model?
  \item \textbf{RQ2 (ablation).} Cơ chế nào của harness tạo ra chất lượng đó, và đóng
  góp của nó có thay đổi theo model scale không?
\end{itemize}

\subsection{Bộ dữ liệu}
\label{sec:dataset}

Chúng tôi dùng ISD-Agent-Bench~\cite{isdagentbench}, với các scenario trải một ma
trận Context~$\times$~ISD. Đánh giá chạy trên split \texttt{test\_90}---90 trong số
các test scenario, phân tầng theo 10 domain $\times$ 3 mức độ khó $\times$ 3 scenario
mỗi ô (seed 42)---cho phủ rộng, cân bằng với chi phí chấm hợp lý.

\subsection{Model và Thiết lập}
\label{sec:setup}

Để thăm dò tính độc lập với model, chúng tôi chạy mọi cấu hình ở hai điểm khác biệt
trên phổ model-scale: \textbf{DeepSeek-v4-flash} (một LLM lớn) và \textbf{Qwen3.5-9B}
(một SLM 9B), với cấu hình giống nhau trừ agent model. Agent model và judge model
được cấu hình riêng để tránh self-preference bias; các judge được giữ cố định qua mọi
lượt chạy và dùng cấu hình multi-judge mặc định của benchmark
(\texttt{openai/gpt-4o-mini} và \texttt{google/gemini-2.5-flash-lite}). Mỗi lượt chạy
ghi một manifest (model, judge, phiên bản prompt, commit, seed, cấu hình decoding)
để tái lập.

\subsection{Các Điều kiện}
\label{sec:conditions}

Với RQ1, AlignmentGraph-ISD được tích hợp làm agent thứ bảy dưới input/output schema
của benchmark và so với sáu agent: Baseline (single prompt), ReAct-ADDIE, ADDIE-Agent,
Dick-Carey-Agent, RPISD-Agent và EduPlanner.

Với RQ2, ablation loại bỏ đúng một cơ chế mỗi lần khỏi harness đầy đủ. Mỗi arm là một
workflow độc lập (không cờ cấu hình dùng chung), nên không cơ chế nào rò rỉ giữa các
arm:
\begin{itemize}
  \item \texttt{full} --- multi-agent decomposition + per-agent verify/repair;
  \item \texttt{wo-ma} --- bỏ decomposition: một agent đơn sinh toàn bộ chương trình
  và tự sửa toàn cục cho chính nó (giữ verify/repair);
  \item \texttt{wo-qc} --- bỏ verify/repair: cùng các specialist thứ tự cố định sinh,
  không có công cụ kiểm soát chất lượng.
\end{itemize}
So \texttt{full} với \texttt{wo-ma} tách riêng multi-agent decomposition; so với
\texttt{wo-qc} tách riêng per-agent verification/repair.

\subsection{Chỉ số}
\label{sec:metrics}

Các chỉ số chính là của benchmark: một ADDIE artifact rubric 100 điểm và một
trajectory score, gộp thành một \emph{Total} composite, mỗi cái là trung bình hai
judge. Chúng tôi bổ sung một \emph{residual alignment-defect audit} tất định:
verifier được chạy lại trên output cuối của mỗi arm và đếm các vi phạm mức nội dung
còn sót. Điều này đo trực tiếp điều kiện (i) của bài toán, nắm bắt một thuộc tính mà
rubric không tưởng thưởng tường minh.

\subsection{Phân tích Thống kê}
\label{sec:stats}

Các so sánh được ghép cặp theo scenario. Chúng tôi dùng kiểm định Wilcoxon
signed-rank (xấp xỉ chuẩn có hiệu chỉnh tie) trên composite score, với hiệu chỉnh
Holm qua các so sánh trong mỗi lượt chạy, và báo cáo hiệu trung bình theo cặp cùng
khoảng tin cậy bootstrap 95\% (10{,}000 resample). Với RQ2, đóng góp của một cơ chế là
$\texttt{full} - \texttt{arm}$, ghép cặp theo scenario.

\subsection{Tái lập}
\label{sec:reproducibility}

ISD-Agent-Bench và sáu baseline của nó là công khai. Mỗi lượt chạy được ghim bằng một
manifest (\cref{sec:setup}); ba ablation arm là các module workflow riêng chọn theo
tên; và các script phân tích (kiểm định ghép cặp và residual audit) được công bố cùng
bài báo.

\section{Kết quả}
\label{sec:results}

Toàn bộ bảng và trị số thống kê trong phần này được sinh tự động từ artifact kết quả
bằng pipeline phân tích (\cref{sec:reproducibility}), nhằm bảo đảm số trong bảng và
trong prose luôn khớp với artifact. \TBD{Các lượt chạy trên dải Qwen đang được thực
hiện; mọi số liệu hiển thị hiện là placeholder ở chế độ demo và sẽ được thay bằng số
thực khi pipeline hoàn tất.}

\subsection{RQ1: Chất lượng theo model scale}
\label{sec:rq1}

\Cref{tab:rq1-ladder} trình bày Total của bảy tác tử trên \texttt{test\_90} theo bốn
kích thước Qwen; mỗi ô là trung bình per-scenario gộp ba lượt chạy độc lập, kèm độ
lệch chuẩn giữa các lượt và $n$ scenario complete-case. Ở kích thước lớn nhất (9B),
AlignmentGraph-ISD đạt Total $\rqOneHarnessTotalNineB$
(SD giữa các lượt $\rqOneHarnessRunSDNineB$), so với $\rqOneBestBaselineTotalNineB$
của đối sánh mạnh nhất tại kích thước này (\rqOneBestBaselineNameNineB); chênh lệch
ghép cặp tương ứng là $+\rqOneDeltaBestNineB$ điểm với khoảng tin cậy bootstrap 95\%
$[\rqOneDeltaBestCILoNineB;\,\rqOneDeltaBestCIHiNineB]$. Ở kích thước nhỏ nhất
(0,8B), harness đạt $\rqOneHarnessTotalZeroEightB$ so với
$\rqOneBestBaselineTotalZeroEightB$ của \rqOneBestBaselineNameZeroEightB{}
(chênh lệch $+\rqOneDeltaBestZeroEightB$;
CI $[\rqOneDeltaBestCILoZeroEightB;\,\rqOneDeltaBestCIHiZeroEightB]$). Trên toàn bộ
24 so sánh ghép cặp (bốn kích thước, sáu đối sánh mỗi kích thước, hiệu chỉnh Holm
trong từng kích thước), $p_{\mathrm{Holm}}$ lớn nhất là $\rqOnePHolmMaxAll$.
\TBD{prose bổ sung: thứ hạng đầy đủ theo từng kích thước, các đối sánh suy giảm mạnh,
và nhận xét về SD giữa các lượt, khi có số thực}.

% AUTO-GENERATED by isd-agent-benchmark/scripts/alignmentgraph-isd-bench/09_gen_paper_tables.py — DO NOT EDIT BY HAND
% tab_rq1_ladder.tex — RQ1 ladder table on ADDIE
% generated_at: 2026-07-30T17:41:41
% source pooled_ladder.json generated_at: 2026-07-30T15:54:18

\begin{table*}[t]
  \caption{RQ1: ISD-Agent-Bench \texttt{test\_90} theo thang kích thước mô hình (Qwen, 3 run độc lập mỗi kích thước). Mỗi ô: trung bình rubric ADDIE per-scenario (/100) — thước đo chất lượng nội dung thiết kế giảng dạy, tức construct mà câu hỏi nghiên cứu nhắm tới, gộp mean-of-runs trên các scenario complete-case; giá trị sau $\pm$ là SD giữa 3 run (run-to-run); $n$ = số scenario đủ cả 3 run. \textbf{Đậm} = tốt nhất theo cột.}
  \label{tab:rq1-ladder}
  \small
  \begin{tabular}{@{}lr@{\hskip 3pt}rr@{\hskip 3pt}rr@{\hskip 3pt}rr@{\hskip 3pt}r@{}}
    \toprule
    & \multicolumn{2}{c}{Qwen3.5-0,8B} & \multicolumn{2}{c}{Qwen3.5-2B} & \multicolumn{2}{c}{Qwen3.5-4B} & \multicolumn{2}{c}{Qwen3.5-9B} \\
    \cmidrule(lr){2-3}\cmidrule(lr){4-5}\cmidrule(lr){6-7}\cmidrule(lr){8-9}
    Agent & ADDIE {\scriptsize$\pm$SD} & $n$ & ADDIE {\scriptsize$\pm$SD} & $n$ & ADDIE {\scriptsize$\pm$SD} & $n$ & ADDIE {\scriptsize$\pm$SD} & $n$ \\
    \midrule
    Baseline & 44{,}52{\scriptsize$\pm$1{,}19} & 5\textsuperscript{*} & 57{,}82{\scriptsize$\pm$2{,}01} & 47\textsuperscript{*} & 78{,}68{\scriptsize$\pm$0{,}48} & 82\textsuperscript{*} & 83{,}45{\scriptsize$\pm$0{,}32} & 86\textsuperscript{*} \\
    EduPlanner & 46{,}61{\scriptsize$\pm$2{,}60} & 7\textsuperscript{*} & 55{,}12{\scriptsize$\pm$0{,}58} & 57\textsuperscript{*} & 77{,}61{\scriptsize$\pm$0{,}27} & 77\textsuperscript{*} & 80{,}31{\scriptsize$\pm$0{,}70} & 85\textsuperscript{*} \\
    ADDIE-Agent & 73{,}55{\scriptsize$\pm$0{,}13} & 90 & 76{,}98{\scriptsize$\pm$0{,}20} & 90 & 79{,}82{\scriptsize$\pm$0{,}05} & 90 & 80{,}42{\scriptsize$\pm$0{,}06} & 90 \\
    RPISD-Agent & 77{,}79{\scriptsize$\pm$0{,}11} & 88\textsuperscript{*} & 79{,}69{\scriptsize$\pm$0{,}20} & 90 & 81{,}03{\scriptsize$\pm$0{,}17} & 90 & 81{,}48{\scriptsize$\pm$0{,}06} & 90 \\
    Dick-Carey-Agent & 79{,}74{\scriptsize$\pm$0{,}31} & 90 & 79{,}62{\scriptsize$\pm$0{,}26} & 90 & 81{,}84{\scriptsize$\pm$0{,}34} & 89\textsuperscript{*} & 78{,}88{\scriptsize$\pm$0{,}13} & 90 \\
    ReAct-ADDIE & 77{,}48{\scriptsize$\pm$0{,}36} & 90 & 82{,}19{\scriptsize$\pm$0{,}25} & 90 & 85{,}08{\scriptsize$\pm$0{,}15} & 90 & \textbf{86{,}30}{\scriptsize$\pm$0{,}08} & 90 \\
    \textbf{AlignmentGraph-ISD} & \textbf{80{,}46}{\scriptsize$\pm$0{,}26} & 90 & \textbf{84{,}69}{\scriptsize$\pm$0{,}02} & 90 & \textbf{85{,}88}{\scriptsize$\pm$0{,}16} & 90 & 86{,}23{\scriptsize$\pm$0{,}15} & 90 \\
    \bottomrule
  \end{tabular}
  \par\smallskip\footnotesize\textsuperscript{*}$n<$ tổng số scenario: có scenario lỗi/thiếu ở ít nhất một run (complete-case; xem stats\_summary.md).
\end{table*}


\paragraph{Interaction theo kích thước mô hình.} Để kiểm tra nhận định về khoảng cách
theo kích thước bằng một phép kiểm định trực tiếp thay vì so hai ước lượng riêng,
chúng tôi báo difference-in-differences của chênh lệch harness trừ Baseline giữa kích
thước lớn nhất và nhỏ nhất: $\rqOneDiDVsBaseline$ điểm (CI 95\%
$[\rqOneDiDVsBaselineCILo;\,\rqOneDiDVsBaselineCIHi]$;
$p=\rqOneDiDVsBaselineP$), cùng độ dốc xu hướng của chênh lệch theo bậc kích thước
($\rqOneTrendSlopeVsBaseline$ điểm/bậc). Kết luận về việc lợi thế của harness có nới
rộng khi mô hình nhỏ đi hay không dựa trên dấu và khoảng tin cậy của các đại lượng
này. \TBD{diễn giải interaction khi có số thực}.

\subsection{RQ2: Constructive alignment theo tiêu chí đã thiết lập}
\label{sec:rq2}

\Cref{tab:rq2-alignment} trình bày các endpoint cấp objective
(\cref{sec:alignment-metric})---full alignment rate (FAR, endpoint chính), assessment
coverage, activity coverage, cognitive congruence---cùng alignment composite cho bảy
tác tử theo bốn kích thước mô hình. Chúng tôi báo cáo ba lớp kết quả: (i) so sánh
ghép cặp giữa harness và từng đối sánh trong từng kích thước trên FAR và composite
(Wilcoxon signed-rank, hiệu chỉnh Holm với family sáu so sánh cho từng chỉ số,
khoảng tin cậy bootstrap); (ii) sự thay đổi của từng endpoint và thành phần khi kích
thước mô hình giảm; và (iii) interaction method~$\times$~model size trên composite
bằng difference-in-differences, kiểm tra trực tiếp giả thuyết rằng khoảng cách
alignment nới rộng khi mô hình nhỏ đi. \TBD{prose kết quả RQ2: mức FAR và composite,
chênh lệch ghép cặp, CI, $p_{\mathrm{Holm}}$ và kết quả interaction, khi có số thực}.

% AUTO-GENERATED by isd-agent-benchmark/scripts/alignmentgraph-isd-bench/09_gen_paper_tables.py — DO NOT EDIT BY HAND
% tab_rq2_alignment.tex — RQ2 objective->assessment similarity x model sizes
% generated_at: 2026-07-30T17:41:41
% source pooled_ladder.json generated_at: 2026-07-30T15:54:18

\begin{table}[t]
  \caption{RQ2: objective$\to$assessment similarity (mean-max rectified cosine, liên tục, không ngưỡng) --- vế \emph{measures} của \textit{aligned(o)}; trung bình per-scenario (SD), gộp mean-of-runs, theo kích thước mô hình. Sáu tín hiệu còn lại của panel: Bảng~\ref{tab:rq2-panel}.}
  \label{tab:rq2-alignment}
  \small
  \begin{tabular}{@{}lrrrr@{}}
    \toprule
    Agent & Qwen3.5-0,8B & Qwen3.5-2B & Qwen3.5-4B & Qwen3.5-9B \\
    \midrule
    Baseline & 0{,}104 (0{,}20) & 0{,}383 (0{,}17) & 0{,}547 (0{,}09) & 0{,}562 (0{,}05) \\
    EduPlanner & 0{,}360 (0{,}23) & 0{,}500 (0{,}16) & 0{,}598 (0{,}09) & 0{,}592 (0{,}07) \\
    ADDIE-Agent & 0{,}595 (0{,}06) & 0{,}606 (0{,}08) & 0{,}608 (0{,}06) & 0{,}586 (0{,}06) \\
    RPISD-Agent & 0{,}510 (0{,}07) & 0{,}519 (0{,}05) & 0{,}466 (0{,}04) & 0{,}451 (0{,}04) \\
    Dick-Carey-Agent & 0{,}441 (0{,}05) & 0{,}499 (0{,}05) & 0{,}536 (0{,}05) & 0{,}557 (0{,}05) \\
    ReAct-ADDIE & 0{,}504 (0{,}11) & 0{,}583 (0{,}10) & 0{,}579 (0{,}06) & 0{,}570 (0{,}06) \\
    \textbf{AlignmentGraph-ISD} & 0{,}851 (0{,}05) & 0{,}850 (0{,}06) & 0{,}806 (0{,}07) & 0{,}832 (0{,}05) \\
    \bottomrule
  \end{tabular}
\end{table}


Vì quy trình đánh giá đọc cùng một văn bản kết xuất cho mọi tác tử và không dùng phán
đoán LLM ở bất kỳ bước tổng hợp nào, kết quả RQ2 là một phép đo tự động của vị từ
$\mathrm{aligned}(o)$ trong \cref{eq:aligned} theo các tiêu chí đã thiết lập, độc lập
với cả verifier của harness lẫn LLM judge của benchmark; các so sánh được diễn giải
ở mức tương đối giữa các tác tử. Mối liên hệ giữa các chỉ số alignment và điểm judge
của RQ1---hai phép đo nhắm vào hai khía cạnh khác nhau của chất lượng---được thảo
luận trong \cref{sec:discussion}.

\section{Thảo luận}
\label{sec:discussion}

\paragraph{Khung diễn giải theo hai câu hỏi nghiên cứu.}
Vì số liệu cuối cùng được điền tự động từ artifact, phần này trình bày khung diễn giải
đã định trước theo từng kịch bản kết quả, thay vì diễn giải hậu nghiệm. Với RQ1, nếu
AlignmentGraph-ISD giữ thứ hạng dẫn đầu trên cả bốn kích thước và
difference-in-differences cho thấy lợi thế lớn hơn ở kích thước nhỏ với khoảng tin cậy
loại trừ 0, kết quả ủng hộ luận điểm harness-over-model~\cite{cho2026harness}: cách tổ
chức môi trường thực thi bù một phần cho giới hạn của mô hình nền, và phần bù này lớn
dần khi mô hình yếu đi. Nếu lợi thế ổn định nhưng không nới rộng, kết luận hẹp hơn:
harness cải thiện chất lượng đồng đều nhưng không đặc biệt hữu ích hơn cho SLM. Nếu
lợi thế thu hẹp ở kích thước nhỏ, giả thuyết thiết kế bị thách thức---các cơ chế của
harness có thể đòi hỏi một năng lực mô hình tối thiểu để vận hành.
\TBD{chốt diễn giải theo kịch bản xảy ra khi có số thực}.

Với RQ2, các chỉ số alignment cho phép tách hai khả năng mà điểm judge không phân biệt
được. Nếu full alignment rate và composite của harness cao hơn các đối sánh và khoảng
cách nới rộng khi mô hình nhỏ đi, điều đó ủng hộ trực tiếp cơ chế trung tâm: alignment
là thuộc tính suy giảm sớm trên mô hình nhỏ, và việc lưu nó thành trạng thái tường
minh giúp giữ lại thuộc tính này. Ngược lại, nếu các chỉ số alignment của harness
không tách khỏi các đối sánh dù điểm judge tách, thì lợi thế của harness đến từ các
khía cạnh chất lượng khác ngoài alignment---một kết quả vẫn hữu ích nhưng đòi hỏi hạ
claim về cơ chế. Sự đối chiếu giữa các chỉ số alignment và điểm judge cũng cho biết
rubric tổng quát nhạy đến đâu với vi phạm alignment.
\TBD{đối chiếu alignment--judge khi có số thực}.

\paragraph{Alignment graph như artifact giải thích được.}
Ngoài vai trò trạng thái cho máy, alignment graph còn là một sản phẩm cho con người.
Toàn bộ node, cạnh và trạng thái kiểm chứng được kết xuất thành file JSON kèm biểu
diễn Graphviz (\cref{sec:aggregation}), nên chuyên gia thiết kế giảng dạy có thể trả
lời các câu hỏi kiểm toán cụ thể---mục tiêu nào chưa có assessment tương ứng, liên kết
nào chưa qua kiểm chứng, sửa đổi nào đã kích hoạt kiểm tra lại---mà không cần đọc lại
toàn bộ văn bản chương trình. Điều này định vị hệ thống là công cụ hỗ trợ ra quyết
định (decision support) trong quy trình educator-in-the-loop: con người giữ quyền rà
soát, chỉnh sửa và override, còn harness bảo đảm rằng các quan hệ cần rà soát là tường
minh và truy vết được. Chúng tôi chủ ý tránh các claim kiểu ``đúng by construction'':
cơ chế chỉ làm cho alignment kiểm tra được và sửa được, không bảo đảm nó đúng.

\paragraph{Triển khai có trách nhiệm.}
Một chương trình sinh tự động có thể sai theo cách khó nhận ra: mục tiêu hoặc
assessment không phù hợp vẫn đọc trôi chảy. Vì vậy, đầu ra của hệ thống cần được
chuyên gia phê duyệt trước khi dùng với người học, đặc biệt trong các ngữ cảnh có hệ
quả cao; artifact đồ thị giảm chi phí của việc rà soát này nhưng không thay thế nó.
Nội dung sinh cũng có thể mang thiên lệch văn hoá, ngôn ngữ hoặc rào cản tiếp cận thừa
hưởng từ mô hình nền; các cơ chế của harness không kiểm tra các thuộc tính này và
chúng cần quy trình rà soát riêng. Cuối cùng, động cơ về riêng tư và tự lưu trữ chỉ
được hiện thực hoá trọn vẹn khi mô hình và judge chạy on-premise; cấu hình thực nghiệm
hiện tại vẫn dùng API bên ngoài (\cref{sec:threats-limitations}).

\section{Nguy cơ đối với tính hợp lệ}
\label{sec:threats-limitations}

\paragraph{Tính hợp lệ nội tại.}
Prompt, mô hình và judge đều có thể ảnh hưởng tới điểm số. Chúng tôi giảm rủi ro này
bằng cách lưu manifest cho từng lượt chạy, giữ cấu hình decoding giống nhau giữa các
tác tử và kích thước, tách mô hình sinh khỏi mô hình judge và sử dụng hai judge. Một
yếu tố gây nhiễu vẫn còn: AlignmentGraph-ISD dùng nhiều lời gọi mô hình hơn Baseline
một-prompt, nên một phần lợi thế trên điểm judge có thể đến từ test-time compute bổ
sung. Nghiên cứu báo số lời gọi và token theo tác tử trong artifact
(\TBD{bảng resource theo tác tử}), nhưng chưa có đối sánh cân bằng ngân sách; kết luận
vì vậy áp dụng cho các tác tử như chúng được thiết kế, không phải cho các phương án đã
cân bằng compute.

\paragraph{Tính hợp lệ khái niệm.}
Rubric ADDIE và điểm trajectory đo chất lượng đầu ra cùng dấu vết thực thi, không đo
kết quả học tập. Thước đo alignment của RQ2 khắc phục việc alignment không được đo
trực tiếp, nhưng mang các giới hạn riêng cần khai báo. Thứ nhất, bộ phân loại Bloom và
encoder có sai số ($\kappa<1$ so với nhãn chuyên gia trong các nghiên cứu
gốc~\cite{waheed2021bloomnet,li2022bloomclassification}); vì cùng một pipeline được áp
dụng cho mọi tác tử, sai số này chủ yếu ảnh hưởng giá trị tuyệt đối, nên kết quả RQ2
được diễn giải là \emph{xếp hạng tương đối giữa các tác tử}, không phải một cuộc kiểm
toán alignment tuyệt đối. Thứ hai, các tài nguyên phân loại Bloom chủ yếu được huấn
luyện trên tiếng Anh; việc chuyển sang scenario tiếng Hàn (qua encoder đa ngữ và
fallback từ điển tiếng Hàn) là một nguy cơ chuyển giao chưa được kiểm chứng độc lập.
Thứ ba, thước đo lượng hoá alignment cấu trúc--ngữ nghĩa giữa các thành phần văn bản;
nó không đo tính đúng đắn sư phạm sâu hơn hay kết quả học tập thực tế. Một đánh giá
với chuyên gia gán nhãn vẫn là bước xác nhận cần thiết.

\paragraph{Tính hợp lệ ngoại tại.}
Các scenario đều là dữ liệu tổng hợp từ một benchmark; kết quả có thể thay đổi trên
brief thực tế, lĩnh vực khác hoặc benchmark ISD khác. Dải kích thước nằm trong một họ
mô hình duy nhất (Qwen) trên cùng runtime: thiết kế này cô lập hiệu ứng kích thước
khỏi khác biệt family, nhưng đổi lại, kết luận về model scale chỉ chắc chắn trong họ
mô hình đã đánh giá và cần được kiểm tra lại trên các họ khác. Ngoài ra, dù động cơ
hướng tới tự lưu trữ, các thí nghiệm gọi mô hình qua API bên ngoài; hành vi trên
hạ tầng on-premise với quantization thực tế chưa được đo.

\paragraph{Tính hợp lệ của kết luận thống kê.}
Các hiệu ứng được ước lượng trên tối đa 90 cặp scenario với ba lượt chạy mỗi tổ hợp;
SD giữa các lượt được báo cáo và các kiểm định dùng trung bình theo scenario, nên
run-to-run variability được lượng hoá nhưng chưa được mô hình hoá đầy đủ (ví dụ bằng
mixed model theo cả scenario và lượt chạy). Complete-case có thể tạo thiên lệch nếu
lỗi không xảy ra ngẫu nhiên; phân tích độ nhạy với scenario lỗi bị phạt điểm
(\cref{sec:stats}) kiểm tra mức nhạy của kết luận với chính sách này. Các scenario
trong cùng lĩnh vực có thể phụ thuộc lẫn nhau; bootstrap theo cụm lĩnh vực là một
kiểm tra bổ sung đáng thực hiện.

\section{Hạn chế và hướng phát triển}
\label{sec:future-work}

Nghiên cứu giới hạn ở việc sinh một chương trình từ một design brief cho trước. Brief
trong thực tế thường hình thành từ trao đổi với các bên liên quan, nên một giao diện hội
thoại để thu thập brief là mở rộng tự nhiên. Harness cũng có thể tiếp nhận tri thức tổ
chức---chẳng hạn rubric nội bộ hoặc quy ước chuyên ngành---qua retrieval, nhưng khả năng
này chưa được đánh giá. Về đo lường, ưu tiên tiếp theo là hiệu chuẩn thước đo alignment
với nhãn của chuyên gia ISD (báo agreement giữa người gán nhãn) và liên hệ các chỉ số
alignment với kết quả học tập thực tế. Về triển khai, một đánh giá trên hạ tầng
on-premise và một nghiên cứu người dùng với chuyên gia thiết kế giảng dạy sử dụng
artifact đồ thị để audit là hai bước cần thiết để kiểm chứng các lợi ích đã nêu trong
động cơ.

\section{Kết luận}
\label{sec:conclusion}

\TBD{}

% ==== BẢN NHÁP CŨ (giữ tham khảo, sẽ viết lại) ====
% Constructive alignment trở nên khó kiểm tra khi objective, hoạt động, assessment và
% bằng chứng evaluation được sinh thành một văn bản trôi chảy duy nhất. Trong serving
% envelope của model nhỏ, harness còn phải duy trì được các quan hệ đó mà không dồn
% toàn bộ tác vụ vào một lời gọi.
%
% Bài báo đưa ra hai lời giải. AlignmentGraph-ISD điều phối các specialist qua một
% graph state tường minh và xuất trạng thái này để audit. Giao thức đánh giá đi kèm
% không dùng LLM và báo cáo một panel phẳng gồm các tín hiệu textual correspondence
% và cognitive demand; giao thức đo các tín hiệu của constructive alignment, không
% phải chính construct đã được validate.
%
% Trên dải Qwen 0.8B--9B, harness sinh output dùng được cho mọi scenario, dẫn panel
% RQ2 theo quy tắc đọc pre-specified ở mọi kích thước, và dẫn ADDIE ở ba kích thước
% nhỏ với lợi thế tương đối lớn nhất về phía đầu nhỏ. Ba tín hiệu judge vẫn đồng
% hạng: \texttt{total\_score} thua một baseline ở kích thước nhỏ nhất, còn ReAct-ISD
% đạt parity ADDIE tại 9B.
%
% Ablation thu hẹp claim về cơ chế. Bỏ graph context kéo panel RQ2 xuống ở mọi kích
% thước trong khi ADDIE gần như bất biến; verifier là repair backstop và nguồn bằng
% chứng audit, không có score effect trong cấu hình shipped. Vì mọi arm đều giữ
% specialist decomposition, đóng góp của decomposition cho robustness chỉ được xác
% lập ở cấp hệ thống.
%
% Tổng thể, bằng chứng ủng hộ việc dùng một graph-guided harness có thể audit trong
% quy trình educator-in-the-loop trên ladder đã đánh giá. Bước tiếp theo là expert
% study blind-coding để kiểm tra agreement máy--người và phân xử hai cách đọc của
% ablation.


\section*{Khả năng tiếp cận dữ liệu}
ISD-Agent-Bench, sáu tác tử đối sánh và pipeline đánh giá đã được công bố công khai
~\cite{isdagentbench}.
Mã nguồn AlignmentGraph-ISD và các artifacts liên quan có thể được truy cập tại [PENDING].

\bibliographystyle{ACM-Reference-Format}
\bibliography{../references}

\end{document}
